\documentclass[11pt,a4paper]{article}

% ──────────────────────────────────────────────────────────────────────────
% AXON-RFC-006 Appendix B (impl) — easynet.web_app
% Status: Draft v0.3 (HTML/EAL dual-canonical: each web_app
%                     carries an EAL source whose hash is in
%                     canonical state; HTML and api_surface are
%                     deterministic projections of the EAL.
%                     Mirrors Appendix A v0.3.)
% Build:  xelatex AXON-RFC-006-B-easynet-webapp.tex
% Companion (normative core):    AXON-RFC-006-stateful-easynet.tex
% Companion (pressure test md):  AXON-RFC-006-stateful-easynet.md (§B)
% Companion (simpler appendix):  AXON-RFC-006-A-easynet-pages.tex
% ──────────────────────────────────────────────────────────────────────────

\usepackage[margin=1in]{geometry}
\usepackage{amsmath,amssymb,amsthm}
\usepackage{xcolor}
\usepackage[colorlinks=true,linkcolor=blue!60!black,citecolor=blue!60!black,urlcolor=blue!60!black]{hyperref}
\usepackage{booktabs}
\usepackage{longtable}
\usepackage{array}
\usepackage{tabularx}
\usepackage{xltabular}
\newcolumntype{L}[1]{>{\raggedright\arraybackslash}p{#1}}
\usepackage{listings}
\usepackage{enumitem}
\usepackage{tcolorbox}
\usepackage{titlesec}
\usepackage{fancyhdr}
\usepackage{parskip}
\usepackage{xspace}

\usepackage{fontspec}
\usepackage{xeCJK}
\setCJKmainfont{PingFang SC}[
    BoldFont={PingFang SC Semibold},
    ItalicFont={PingFang SC Light}
]

\pagestyle{fancy}
\fancyhf{}
\fancyhead[L]{\small AXON-RFC-006 Appendix B}
\fancyhead[R]{\small easynet.web\_app --- Technical Design}
\fancyfoot[C]{\thepage}
\renewcommand{\headrulewidth}{0.4pt}

\titleformat{\section}{\Large\bfseries\color{black!85}}{\thesection}{0.7em}{}
\titleformat{\subsection}{\large\bfseries\color{black!75}}{\thesubsection}{0.6em}{}
\titleformat{\subsubsection}{\normalsize\bfseries\color{black!70}}{\thesubsubsection}{0.5em}{}

\definecolor{normativecolor}{RGB}{180,30,30}
\definecolor{derivedcolor}{RGB}{30,90,180}
\definecolor{deferredcolor}{RGB}{120,90,30}
\definecolor{codebg}{RGB}{245,245,248}
\definecolor{codeborder}{RGB}{210,210,220}

\newcommand{\normsv}{\textcolor{normativecolor}{\textbf{[normative]}}\xspace}
\newcommand{\derived}{\textcolor{derivedcolor}{\textbf{[derived]}}\xspace}
\newcommand{\deferred}{\textcolor{deferredcolor}{\textbf{[deferred]}}\xspace}

\lstdefinestyle{ealstyle}{
    backgroundcolor=\color{codebg},
    basicstyle=\ttfamily\small,
    frame=single,
    rulecolor=\color{codeborder},
    breaklines=true,
    showstringspaces=false,
    keywordstyle=\color{blue!70!black}\bfseries,
    commentstyle=\color{green!40!black}\itshape,
    stringstyle=\color{red!60!black},
    morekeywords={state_type,state_key,canonical,operational,transition,receipt,view,principal,agent,owner_signature,attempt_id,transition_id,routes_manifest,api_abilities,actions,fn,let,struct,impl,pub,async,await,location,server,proxy_pass,call,on,with,timeout,mission,inputs,return},
    columns=fullflexible,
    keepspaces=true,
    aboveskip=0.6em,
    belowskip=0.6em,
}
\lstset{style=ealstyle}

\newtheorem{theorem}{Theorem}
\newtheorem{lemma}[theorem]{Lemma}
\newtheorem{corollary}[theorem]{Corollary}
\theoremstyle{definition}
\newtheorem{definition}{Definition}

\newtcolorbox{invariant}[1][Invariant]{
    colback=red!4!white,
    colframe=red!50!black,
    fonttitle=\bfseries,
    title={Invariant --- #1},
    sharp corners=south,
    boxrule=0.6pt,
}
\newtcolorbox{notebox}[1][]{
    colback=blue!4!white,
    colframe=blue!50!black,
    fonttitle=\bfseries,
    title={Note},
    sharp corners=south,
    boxrule=0.6pt,
    #1
}
\newtcolorbox{warningbox}[1][]{
    colback=orange!6!white,
    colframe=orange!70!black,
    fonttitle=\bfseries,
    title={Transitional misalignment},
    sharp corners=south,
    boxrule=0.6pt,
    #1
}
\newtcolorbox{decisionbox}[1][]{
    colback=purple!4!white,
    colframe=purple!60!black,
    fonttitle=\bfseries,
    title={Open decision},
    sharp corners=south,
    boxrule=0.6pt,
    #1
}

\title{
    \vspace{-2em}
    \textbf{\Large AXON-RFC-006 Appendix B: easynet.web\_app}\\[0.3em]
    \large Technical design, build pipeline, Hub routing,
    \texttt{actions.*} namespace, and implementation plan
    for full-stack agent-deployed applications
}
\author{
    Silan Hu \\ \small National University of Singapore
}
\date{Draft v0.3, \today}

\begin{document}
\maketitle

\begin{abstract}
\noindent
This appendix specifies \texttt{easynet.web\_app}, the
state\_type that supports a full modern web deployment:
a real React/Vue/Next-style frontend project authored by an
agent (typically Claude Code) inside a project directory under
\texttt{\textasciitilde/.easynet/pages-project/<slug>/}, built into
a static \texttt{dist/} bundle, and served alongside a
backend whose endpoints are not pretend HTTP handlers but
genuine EasyNet abilities. The Hub plays the role of a
state-aware nginx: it serves the frontend bundle, translates
\texttt{POST /api/<verb>} requests into ability invocations on
the owner agent, proxies WebSocket subscriptions to ability
streams, and exposes a structured \texttt{agent\_view} on
\texttt{/\_\_easynet/agent}. Beyond serve-and-translate, this
appendix introduces the \texttt{actions.*} ability namespace ---
queryable endpoints (\texttt{actions.list\_for\_page},
\texttt{actions.get}, \texttt{actions.search}) through which a
peer agent obtains EAL fragments and missions \emph{relevant
to a specific page or intent}, on demand, without bulk-downloading
a static manifest. The agent-skill seeder
(\texttt{easynet-pages-author}) installs in every agent's
workspace at first daemon contact so that authoring a web\_app is
a learnable skill, not undocumented protocol knowledge. The
appendix covers: project location and lifecycle; the
state machine (canonical + runtime overlay); the long-running
build pipeline; the Hub routing matrix; the \texttt{actions.*}
contract and the \texttt{/easynet-for-agents/<verb>} HTTP
mirror; agent\_view's role in this redesign; the five PR slices
required to land it in \texttt{EasyNet-Cli}; and an end-to-end
worked example of an agent authoring a Vite + React shop,
declaring its API and actions surface, building, publishing,
serving, and being audited.
\end{abstract}

\tableofcontents
\vspace{1em}
\hrule
\vspace{1em}

% ──────────────────────────────────────────────────────────────────────────
\section{Position and lineage}
\label{sec:position}

\derived{} The web evolved from a static-document medium into a
full application platform; this appendix tracks the same arc as
the second of two appendices in RFC-006:

\begin{center}
\renewcommand{\arraystretch}{1.25}
\begin{tabularx}{\linewidth}{l X}
\toprule
\textbf{Appendix} & \textbf{What it models} \\
\midrule
A (\texttt{easynet.page})    & A single document. Author writes HTML / Markdown; Hub serves it. The minimal reference: one canonical state, no runtime overlay, no build step, no API surface. \\
B (\texttt{easynet.web\_app})& A full application. Author writes a real frontend project, a build pipeline produces a static bundle, the application's API endpoints ARE EasyNet abilities, and a peer agent can query the owner for ``what can I do on this page?'' through the \texttt{actions.*} ability namespace. \\
\bottomrule
\end{tabularx}
\end{center}

\noindent We keep both. Pages remains the simplest verification
of the protocol's normative core; web\_app is the first state
type that exercises long-running canonical transitions, runtime
overlay, full Hub routing, and agent-actionable surface
discovery.

\subsection{Why this is not nginx}

Nginx serves a static directory and reverse-proxies to a
separately-deployed backend. The backend is a black box ---
nginx neither understands its state, nor knows what endpoints
it offers, nor can answer ``what can a caller do on this page''
to a peer asking through structured discovery. The Hub
specified here serves the same kinds of bytes nginx would, but
its routing decisions are derived from the state object's
canonical state, its API endpoints are abilities introspectable
through the same registry that backs every other EasyNet caller,
and its discovery surface (\texttt{actions.*}) is not a
side-band manifest but a first-class queryable namespace owned
by the owner agent. Nginx is HTTP plumbing; the Hub is a
state-projection layer that happens to speak HTTP at its outer
boundary, with an EAL-aware dual on the inside.

\subsection{What changed from v0.1}

Three structural changes since the previous draft:

\begin{enumerate}[leftmargin=2em,itemsep=0.18em]
    \item \textbf{Project home.} Project directory moved from
    \texttt{<workspace>/web\_apps/<slug>/} to a dedicated, global
    \texttt{\textasciitilde/.easynet/pages-project/<slug>/}.
    Rationale: the project belongs to a state object whose
    owner is already encoded in the state\_key, not to the
    agent's home directory.

    \item \textbf{Agent skill seeded by default.} A bundled
    skill, \texttt{easynet-pages-author}, is installed into every
    agent's workspace at first daemon contact, mirroring the
    way \texttt{easynet-ability-crud} is already seeded. The
    skill teaches an agent how to scaffold, declare, build,
    serve, and discover web\_apps. Authoring is a learnable
    skill, not folklore.

    \item \textbf{actions as a queryable ability namespace.} The
    previous draft proposed an ``actions\_view'' that was a
    static directory of EAL files served alongside the page.
    That design forced peer agents to bulk-download every
    action whether or not relevant to their current page or
    intent --- which is the opposite of how a peer would want
    to consume agent-actionable surface. The new design
    introduces \texttt{actions.*} as an EasyNet ability
    namespace: a peer agent calls
    \texttt{<owner>.actions.list\_for\_page(page\_path)} and
    receives only the actions relevant to that page, on demand.
    The HTTP mirror at \texttt{/easynet-for-agents/<verb>} exists
    for browser-side use but the protocol-level contract is the
    ability, not the URL.
\end{enumerate}

% ──────────────────────────────────────────────────────────────────────────
\section{Goals and non-goals}
\label{sec:goals}

\subsection{Goals}

\normsv{} An implementation conforming to this appendix:

\begin{enumerate}[leftmargin=2.2em,itemsep=0.18em]
    \item Lets an agent author a real frontend project under
    \texttt{\textasciitilde/.easynet/pages-project/<slug>/} using
    whatever toolchain the agent prefers (Vite, Next.js, Remix,
    \ldots); EasyNet does not invent a project structure.
    \item Provides a long-running canonical transition,
    \texttt{webapp.build}, that runs the project's declared
    build command, captures the resulting \texttt{dist/} as an
    immutable content-addressed bundle, and advances canonical
    state on success or failure.
    \item Provides runtime overlay transitions
    (\texttt{webapp.serve}, \texttt{webapp.stop}) that bind
    a Hub mapping; canonical state is unchanged.
    \item Translates HTTP requests on the Hub along three
    orthogonal channels: static-bundle fetch, API ability
    invocation (\texttt{/api/<verb>}), and WebSocket ability
    subscription (\texttt{/ws/<topic>}).
    \item Exposes \texttt{agent\_view} on
    \texttt{/\_\_easynet/agent} as a typed JSON document of the
    application's state and discoverability metadata, including
    a pointer to the \texttt{actions.*} namespace.
    \item Defines the \texttt{actions.*} ability namespace
    (\S\ref{sec:actions}): a queryable, owner-implemented
    surface through which peer agents request EAL fragments
    and missions on demand. The HTTP mirror
    \texttt{/easynet-for-agents/<verb>} is an alias for ability
    invocation, not a static directory.
    \item Seeds the \texttt{easynet-pages-author} skill into
    each agent's workspace so authoring a web\_app is a
    documented capability of the agent itself
    (\S\ref{sec:skill-seed}).
    \item Is implementable in five PR slices
    (\S\ref{sec:impl-plan}), each independently reviewable.
\end{enumerate}

\subsection{Non-goals}

\derived{} Out of scope for this appendix:

\begin{itemize}[leftmargin=2.2em,itemsep=0.18em]
    \item TLS termination, HTTP/2 negotiation, multi-instance
    Hub clusters --- operations work, not protocol work.
    \item Choice of frontend framework --- the appendix uses
    React/Vite as the worked example because Claude Code
    defaults to it; the protocol is framework-agnostic.
    \item Server-side rendering with persistent process ---
    Phase 1 is static + Hub-routed API + WebSocket; SSR
    requiring a Node.js process per request is RFC-006-B v2.
    \item Cross-realm or cross-node application federation ---
    Phase 1 is single owner / single node.
    \item Authenticated humans (\texttt{principal/human-auth})
    --- Phase 1 is \texttt{principal/human-anon} only.
    \item ``Three-view equivalence'' as a strict protocol
    invariant --- enforcing that \texttt{human\_view} and
    actions content carry exactly the same information is
    desirable but not protocol-encodable in v1; the
    \texttt{actions.*} namespace is the owner's voluntary
    statement of agent-actionable surface, not a derivation.
\end{itemize}

% ──────────────────────────────────────────────────────────────────────────
\section{Project model}
\label{sec:project}

\subsection{Project location}

\normsv{} The state object's source-of-truth on disk is a
project directory at:

\begin{lstlisting}
~/.easynet/pages-project/<slug>/
├── easynet.app.toml           # the EasyNet-side declaration
├── package.json               # owned by the agent's toolchain
├── package-lock.json          # ditto
├── src/                       # framework-shaped source
├── public/
├── vite.config.ts             # or next.config.js, etc.
└── (optionally) anything else the framework needs
\end{lstlisting}

EasyNet \textbf{owns} only \texttt{easynet.app.toml}; the rest is
the framework's. The location is global (one directory tree per
host, slug-keyed) rather than per-agent (\texttt{<workspace>/...})
because the state\_key already carries owner identity ---
nesting the directory under the agent's workspace would encode
ownership twice and complicate Hub access (the Hub reads
artifacts via the blob store, but a debug operator inspecting a
project should not have to know which agent owns a given slug).

The directory is created by the daemon on
\texttt{webapp.create}; the agent populates it with framework
files via its own toolchain (\texttt{npm create vite}, etc.).
The daemon does not write into this directory after creation
except through transitions on the state object's behalf
(e.g. clearing on \texttt{webapp.delete}). The blob store at
\texttt{\textasciitilde/.easynet/blobs/} is the canonical source
of bytes served by the Hub; the project directory is the
agent's working copy.

\begin{invariant}[WB-INV-1]
After a successful \texttt{webapp.build}, the Hub serves bytes
exclusively from the blob store keyed by
\texttt{build\_artifact\_hash}. The project directory's
\texttt{dist/} may exist on disk for the agent's debugging
convenience but is \textbf{never} read by the Hub. Reading the
project's \texttt{dist/} for serving purposes is non-conformant.
\end{invariant}

\subsection{easynet.app.toml}

\normsv{} The single file EasyNet reads to understand a web\_app:

\begin{lstlisting}
# easynet.app.toml — the only file EasyNet authors and reads.
# The rest of the project belongs to the framework.

[manifest]
name         = "shop"
description  = "Tiny ceramic shop, agent-built."
visibility   = "public"            # private | realm | public

[build]
command      = "npm install && npm run build"
dist_dir     = "dist"              # relative to project root
build_timeout_secs = 600
node_version = "20"                # advisory; executor may pin

[serve]
mode         = "static-only"
                # "static-only" : no backend process; API is satisfied
                #                 by ability invocations from the Hub.
                #                 The default and the Phase-1
                #                 production shape.
                # "process"      : a backend process is started;
                #                 reserved for v2 (SSR support).

[api]
# Abilities the FRONTEND is allowed to call through the Hub's
# /api/<verb> route. Each entry is a fully-qualified ability name
# the owner agent has registered. The Hub will refuse to translate
# a /api/<verb> request whose verb is not on this list.
abilities = [
  "shop.list_items",
  "shop.checkout",
  "shop.order_status",
]

[actions]
# Owner-implemented `actions.*` abilities exposed via the
# /easynet-for-agents/<verb> mirror. Each entry is a verb the
# owner agent has registered under the actions.* namespace.
# Peer agents call these to discover EAL fragments and missions
# on demand. See §6 for the protocol.
abilities = [
  "actions.list_for_page",
  "actions.get",
  "actions.search",
]

[routes]
# nginx-shaped routing rules. Order is significant.
rules = [
  { match = "exact:/",                      action = "static:dist/index.html" },
  { match = "prefix:/assets/",              action = "static:dist/assets/" },
  { match = "regex:^/api/([a-z][a-z0-9_]*)$",
    action = "api:$1",
    methods = ["GET","POST"] },
  { match = "regex:^/easynet-for-agents/([a-z][a-z0-9_]*)$",
    action = "actions:$1",
    methods = ["GET","POST"] },
  { match = "prefix:/ws/",                  action = "ws:$path" },
  { match = "exact:/__easynet/agent",       action = "agent_view" },
  { match = "exact:/__easynet/health",      action = "health" },
  { match = "exact:/__easynet/receipts",    action = "receipts" },
  { match = "fallback",                     action = "spa:dist/index.html" },
]

[csp]
extra_script_src  = []
extra_style_src   = ["'unsafe-inline'"]
extra_connect_src = ["self"]
\end{lstlisting}

A minimal app gets by with three sections (\texttt{[manifest]},
\texttt{[build]}, \texttt{[api]}); \texttt{[actions]} is optional
(\S\ref{sec:actions} explains why) and routes default to the
standard SPA + actions layout shown above.

\subsection{Why declared API and actions surfaces, not auto-detected}

\normsv{} Both the \texttt{[api]} and \texttt{[actions]} surfaces
are declared, not discovered. The Hub refuses to translate
\texttt{/api/<verb>} or \texttt{/easynet-for-agents/<verb>}
when the corresponding verb is not declared.

\begin{invariant}[WB-INV-2]
The Hub \textbf{MUST} refuse to translate \texttt{/api/<verb>}
unless \texttt{<verb>} appears in the web\_app's currently-Built
canonical state's \texttt{api\_surface\_hash} (the canonicalised
\texttt{[api].abilities} list). The same rule applies to
\texttt{/easynet-for-agents/<verb>} relative to
\texttt{actions\_surface\_hash}. The refusal MUST be HTTP 404
(not 403) to avoid confirming the existence of internal
abilities --- the application-layer instance of the
trust-boundary discipline TR-INV-12 establishes at the identity
layer.
\end{invariant}

\subsection{Agent skill seeded by default}
\label{sec:skill-seed}

\normsv{} The daemon, on the first contact with each agent,
seeds a bundled skill named \texttt{easynet-pages-author} into
that agent's workspace at the standard skill location:

\begin{lstlisting}
<workspace>/skills/easynet-pages-author/SKILL.md
<workspace>/.agents/skills/easynet-pages-author/SKILL.md
\end{lstlisting}

The skill is identical in shape to the existing
\texttt{easynet-ability-crud} skill (commit 3bc6c70). It
documents:

\begin{itemize}[leftmargin=2em,itemsep=0.15em]
    \item how to scaffold a frontend project
    (\texttt{npm create vite@latest <slug> -- --template react-ts});
    \item the standard project location
    (\texttt{\textasciitilde/.easynet/pages-project/<slug>/}) ---
    to be used \emph{after} \texttt{webapp.create} mints the
    state\_key and creates the directory; the agent does
    \texttt{cd} into the directory the daemon prepares;
    \item the structure of \texttt{easynet.app.toml}, including
    the \texttt{[api]} and \texttt{[actions]} sections;
    \item how to register \texttt{api.*} abilities through the
    existing \texttt{easynet-ability-crud} skill;
    \item how to implement the \texttt{actions.*} namespace
    (\S\ref{sec:actions}) so that peer agents can discover the
    application's actionable surface;
    \item the lifecycle: \texttt{webapp.create} $\to$
    populate project $\to$ \texttt{webapp.build} $\to$
    \texttt{webapp.serve};
    \item how to invoke \texttt{webapp.update\_config} for
    routing or CSP changes that don't require a rebuild;
    \item troubleshooting (build failures, port conflicts,
    visibility-pinned 404s).
\end{itemize}

The skill is a Phase-1 deliverable in P-B5 (see
\S\ref{sec:impl-plan}). It is not part of canonical state ---
it lives in the agent's workspace, owned by the agent, and is
re-seeded only when the daemon detects a missing or
out-of-date skill at boot.

% ──────────────────────────────────────────────────────────────────────────
\section{State machine}
\label{sec:state-machine}

\subsection{state\_type and state\_key}

\normsv{}

\begin{lstlisting}
state_type:  easynet.web_app

state_key:
    (realm: string, owner_agent_uri: string, app_slug: string)

URA form:
  easynet:///r/<realm>/agent/<owner>/web_app/<slug>
\end{lstlisting}

\texttt{app\_slug} matches the slug ABNF defined in
Appendix A's pages section. The directory at
\texttt{\textasciitilde/.easynet/pages-project/<app\_slug>/} is
created on \texttt{webapp.create} and removed on
\texttt{webapp.delete}.

\subsection{Canonical fields}

\normsv{} Per CSH-INV-2 (main \S2.5) every large value is held
by content hash:

\begin{lstlisting}
canonical_fields(web_app) = {
    manifest_hash:        bytes(32),  // [manifest] section
    build_config_hash:    bytes(32),  // [build] section
    eal_hash:             bytes(32),  // SHA-256 of the web_app's
                                      // EAL source (markdown
                                      // documentation in `// ...`
                                      // comments + mission decls).
                                      // Per Appendix A PA-INV-2,
                                      // EAL is the canonical
                                      // agent-facing artifact;
                                      // api_surface_hash is derived
                                      // from EAL at update_config
                                      // time, not separately authored.
    api_surface_hash:     bytes(32),  // [api].abilities (sorted);
                                      // recomputed from EAL by the
                                      // daemon, not written by hand
    actions_surface_hash: bytes(32),  // [actions].abilities (sorted)
    routes_manifest_hash: bytes(32),  // [routes].rules (in order)
    csp_hash:             bytes(32),  // [csp] section
    source_snapshot_hash: bytes(32),  // last successful build's
                                      // source tree hash (null
                                      // before first build)
    build_artifact_hash:  bytes(32),  // dist tree hash (null
                                      // until first Built)
    state_value:          string,
    visibility:           string,
    version:              uint64,
}
\end{lstlisting}

Every field is tagged \texttt{canonical:true} in the schema;
schema-driven projection $\pi$ produces the projection over
exactly these keys.

\subsection{Canonical state space}

\begin{center}
\renewcommand{\arraystretch}{1.25}
\begin{tabularx}{\linewidth}{l X}
\toprule
\textbf{state\_value} & \textbf{Meaning} \\
\midrule
\texttt{Created}   & easynet.app.toml registered; project directory exists; no build attempted yet. \\
\texttt{Building}  & A \texttt{webapp.build} is in flight. Transient. Visible only via OperationalReceipts; never the terminal canonical state. \\
\texttt{Built}     & At least one successful build exists. \texttt{build\_artifact\_hash} non-null. May be Served. \\
\texttt{Failed}    & The most recent build failed. Owner can iterate and retry. \\
\texttt{Deleted}   & Terminal; project directory and blob references cleared. \\
\bottomrule
\end{tabularx}
\end{center}

\subsection{Runtime overlay state space}

\begin{center}
\renewcommand{\arraystretch}{1.25}
\begin{tabularx}{\linewidth}{l X}
\toprule
\textbf{runtime} & \textbf{Meaning} \\
\midrule
\texttt{Stopped}  & No serving overlay associated with this app. \\
\texttt{Starting} & \texttt{webapp.serve} invoked; Hub is binding the slug into its router. \\
\texttt{Running}  & Hub is actively translating \texttt{/<slug>/...} requests for this app. \\
\texttt{Crashed}  & The Hub lost its mapping (process restart, config-reload error, etc.); canonical state untouched. Recoverable via \texttt{webapp.serve}. \\
\bottomrule
\end{tabularx}
\end{center}

\subsection{Transition list}

\begin{center}
\renewcommand{\arraystretch}{1.2}
\begin{tabularx}{\linewidth}{l l X}
\toprule
\textbf{Transition} & \textbf{Class} & \textbf{Notes} \\
\midrule
\texttt{webapp.create@v1}      & canonical   & $\emptyset \to$ Created. Mints state\_key, creates project directory, persists hashes of declared sections. \\
\texttt{webapp.update\_config@v1}
                               & canonical   & Created/Built/Failed $\to$ same. Replaces declared config (manifest, build, api, actions, routes, csp). Does NOT trigger a build. \\
\texttt{webapp.build@v1}       & canonical   & Created/Built/Failed $\to$ Building (transient) $\to$ Built|Failed. Long-running. See \S\ref{sec:build-pipeline}. \\
\texttt{webapp.delete@v1}      & canonical   & * $\to$ Deleted. Implicit \texttt{webapp.stop} as side effect; clears project directory. \\
\texttt{webapp.serve@v1}       & operational & Stopped/Crashed $\to$ Starting $\to$ Running. Requires canonical = Built. \\
\texttt{webapp.stop@v1}        & operational & Starting/Running/Crashed $\to$ Stopped. \\
\texttt{webapp.get@v1}         & query       & Returns canonical + runtime snapshot. \\
\texttt{webapp.get\_eal@v1}    & query       & Mandatory. Returns the UTF-8 EAL source whose SHA-256 equals \texttt{eal\_hash}. HTTP mirror: \texttt{GET /\_\_easynet/eal}. Mirrors Appendix A PA-INV-2. \\
\texttt{webapp.list@v1}        & query       & Lists web\_apps owned by an agent. \\
\bottomrule
\end{tabularx}
\end{center}

% ──────────────────────────────────────────────────────────────────────────
\section{Build pipeline}
\label{sec:build-pipeline}

\subsection{Steps of webapp.build@v1}

\normsv{} On \texttt{webapp.build} the executor:

\begin{enumerate}[leftmargin=2.2em,itemsep=0.18em]
    \item Reads
    \texttt{\textasciitilde/.easynet/pages-project/<slug>/easynet.app.toml}
    and validates it.
    \item Computes \texttt{source\_snapshot\_hash}: a Merkle hash
    over the project directory excluding \texttt{node\_modules/},
    \texttt{dist/}, \texttt{.git/}, and any path the framework's
    \texttt{.gitignore} excludes.
    \item Emits \texttt{build.started} OperationalReceipt with
    the source-snapshot hash, declared command, expected timeout.
    \item Spawns the build command in the project directory.
    Stdout / stderr are tee'd; periodic \texttt{build.progress}
    OperationalReceipts are emitted (lossy).
    \item On exit:
    \begin{itemize}[leftmargin=1.5em]
        \item Status 0 + \texttt{dist\_dir} non-empty: walk the
        \texttt{dist\_dir} tree; compute
        \texttt{build\_artifact\_hash} as a Merkle hash;
        write the entire tree to the blob store as a TreeBlob
        (a JCS-canonical manifest of file paths and individual
        blob hashes; the file blobs are stored as ordinary
        content-addressed blobs).
        \item Status non-zero or empty \texttt{dist\_dir}:
        prepare a failure-reason payload and store its hash.
    \end{itemize}
    \item Emit one CanonicalReceipt:
    \texttt{build.completed} (success) or \texttt{build.failed}
    (failure). The sole canonical advance for the transition;
    \texttt{post\_state\_hash} reflects Built or Failed;
    \texttt{state\_version} +1.
\end{enumerate}

The transient \texttt{Building} state value, the streaming
progress, and the build-tool stdout are all observability;
they never enter the canonical hash.

\subsection{TreeBlob format}

\normsv{}

\begin{lstlisting}
TreeBlob (JCS-canonical JSON)
{
  "version": 1,
  "entries": [
    { "path": "index.html",
      "size": 2148,
      "blob": "<sha256-of-file-bytes>",
      "mode": "0644" },
    { "path": "assets/index-DfX9k.js",
      "size": 187321,
      "blob": "<sha256>",
      "mode": "0644" },
    ...
  ]
}
\end{lstlisting}

The tree's hash is \texttt{sha256(jcs(TreeBlob))} ---
\texttt{build\_artifact\_hash}. Each individual file is also a
content-addressed blob.

\subsection{Build executor}

\normsv{} Phase 1 ships a local build executor: the build runs
in a subprocess of the EasyNet daemon owning the agent. The
owner signs the resulting CanonicalReceipt directly. Delegated
build (source on agent A, build on agent B) is the future:
the protocol shape is already in place (main \S5.3) but the
executor is not.

\begin{decisionbox}
\textbf{Open: build sandboxing.} The build subprocess inherits
daemon process privileges. Phase 1 mitigation: build runs with
a restricted \texttt{PATH}, no inherited env vars except
\texttt{NODE\_VERSION} and \texttt{HOME}, stricter
\texttt{ulimit}. Full sandboxing (seccomp / pledge / Docker)
is RFC-006-B v2.
\end{decisionbox}

% ──────────────────────────────────────────────────────────────────────────
\section{The actions.* ability namespace}
\label{sec:actions}

\subsection{Why a namespace, not a static directory}

\derived{} A peer agent visiting a web\_app does not want to
download every EAL fragment the owner ever wrote; it wants
\emph{the actions relevant to the page or intent it has now}.
Static directories of EAL files force bulk download and
client-side filtering --- exactly the inefficiency the previous
draft had. The redesign treats the agent-actionable surface as
a queryable resource.

The owner agent registers abilities under the \texttt{actions.*}
namespace. A peer agent --- or the Hub on behalf of a browser
--- invokes those abilities by name. The owner's implementation
decides what to return: the page-relevant set, an intent-matched
set, the full directory, a single named action, or any future
form. Because actions are abilities, every existing protocol
mechanism applies: \texttt{ability\_surface}, schema, scope
check, receipts, audit, federation routing.

\subsection{Recommended verbs}

\normsv{} Owners are encouraged but not required to implement a
standard set:

\begin{center}
\renewcommand{\arraystretch}{1.25}
\begin{tabularx}{\linewidth}{l X}
\toprule
\textbf{Ability} & \textbf{Purpose} \\
\midrule
\texttt{actions.list\_for\_page(page\_path)}
    & Return the actions relevant to \texttt{page\_path}. The most
      common entry point: an agent that just landed on a page
      asks ``what can I do here.'' \\
\texttt{actions.get(name)}
    & Return the full definition of one named action ---
      typically the EAL or mission text plus a structured
      \texttt{inputs} schema. \\
\texttt{actions.search(intent)}
    & Return a ranked list of action descriptors matching a
      free-text intent string. Owner-implementation-defined
      ranking. \\
\texttt{actions.list\_all()}
    & Return the full directory. Equivalent to a sitemap; useful
      for crawlers and for mirroring. \\
\bottomrule
\end{tabularx}
\end{center}

The verb set is a \emph{convention}, not protocol. An owner may
implement a subset, a superset, or none; an agent invoking a
non-existent verb gets a normal \texttt{MethodNotFound} error.
The Hub's \texttt{/easynet-for-agents/<verb>} mirror simply
forwards to whatever verb the owner registered.

\subsection{Action descriptors}

\normsv{} A descriptor returned by any \texttt{actions.*} verb
has the shape:

\begin{lstlisting}
{
  "name":     "checkout",          // owner-namespaced
  "kind":     "eal" | "mission" | "form_map",
  "summary":  "Place a single order.",
  "url":      "/easynet-for-agents/get?name=checkout",
                                   // dereference (POST with params)
                                   // to receive a specialised mission
  "params_schema": {               // schema for the params object
                                   //   passed to actions.get;
                                   //   owner inlines them into the
                                   //   returned mission as literals
    "item_id": { "type":"string", "required":true,
                 "doc":"see actions.search to find IDs" },
    "qty":     { "type":"uint",   "default":1 }
  },
  "result_render": {               // optional, for form_map
    "success_template": "Order {order_id}, ETA {eta}",
    "error_template":   "Sorry, {error_message}"
  }
}
\end{lstlisting}

The body of an action --- an EAL fragment, a mission, or a
form\_map JSON --- is whatever \texttt{actions.get} returns.
Because the body is a value returned by an ability, not a file
on disk, the owner may compute it dynamically (e.g. inject the
current user's session into a mission) without violating any
protocol rule. View determinism (TR-INV-6) does not apply
because actions are abilities, not views.

\subsection{What EAL actually is (and what it isn't)}
\label{sec:eal-grammar-note}

\normsv{} The mission bodies that \texttt{actions.get} returns
are EAL programs in the syntax this codebase already implements
(\texttt{src/eal/}). EAL is a deliberately-narrow language:

\begin{itemize}[leftmargin=2em,itemsep=0.15em]
    \item A mission is a \emph{static DAG} of \texttt{call} steps.
    No \texttt{if}, no \texttt{fail}, no expression evaluation,
    no \texttt{return}. Every step's args are filled in when the
    mission text is written.
    \item Argument syntax is \texttt{with \{ key = value, \ldots \}}
    (\texttt{=}, not \texttt{:}). Values are string / int / float /
    bool literals, or \texttt{prev\_step.output} variable references.
    \item Bindings: \texttt{let x = call \ldots}. Variables can be
    referenced only as \texttt{x.output}.
    \item Step options are exactly \texttt{timeout N},
    \texttt{retries N}, \texttt{on\_failure \{abort | skip | retry
    | continue\}}, \texttt{optional}.
    \item The only block form is \texttt{loop} with
    \texttt{body \{\ldots\} verify \{\ldots\}}.
\end{itemize}

This is the actual grammar in \texttt{src/eal/lexer.rs} and
\texttt{src/eal/ast.rs}. There is no \texttt{inputs} clause, no
\texttt{\$variable} interpolation, no in-mission control flow
beyond the \texttt{loop} block.

\textbf{Implication for actions.} Because mission text is
static, an action that needs caller-supplied parameters cannot
``accept inputs'' in the body. Instead the protocol uses
\textbf{owner-side specialisation}: the peer agent passes the
parameters it wants to \texttt{actions.get(name, params)};
the owner agent inlines those parameters as literals into the
mission text it returns. The peer then runs the resulting,
fully-literal mission verbatim through \texttt{mission run}.

This preserves the receipt chain. Each \texttt{call} in the
mission is a real EasyNet invocation that goes through the
existing dispatcher and produces a CanonicalReceipt; the
mission's first step cites \texttt{actions.get}'s receipt as
its \texttt{causal\_context.Scalar} so a later auditor can
prove ``this mission run originated from a specific
\texttt{actions.get} call to alice.''

\subsection{Example: actions.list\_for\_page + actions.get}

\begin{lstlisting}
$ easynet ability invoke alice.actions.list_for_page \
    --args '{"page_path":"/cart"}'

[
  { "name":"checkout", "kind":"mission",
    "summary":"Buy the items currently in this user's cart.",
    "url":"/easynet-for-agents/get?name=checkout",
    "params_schema":{                        // peer passes these to
      "address":{"type":"string","required":true},  // actions.get
      "qty":    {"type":"uint",  "default":1}
    } },
  { "name":"clear_cart", "kind":"eal",
    "summary":"Empty the cart.",
    "url":"/easynet-for-agents/get?name=clear_cart" }
]

# Peer asks owner to specialise the checkout mission for its
# concrete inputs. Owner inlines them as literals into a real EAL
# mission (no $variable, no inputs block).
$ easynet ability invoke alice.actions.get \
    --args '{"name":"checkout",
             "params":{"address":"123 Main St","qty":1}}'

mission "checkout" {
  let cart = call "shop.cart_contents" on "alice" with {} timeout 5

  let order = call "shop.checkout" on "alice" with {
    items   = cart.output,
    address = "123 Main St",
    qty     = 1
  } timeout 30
}
\end{lstlisting}

The peer agent now runs the returned text directly:

\begin{lstlisting}
$ easynet mission run /tmp/checkout.eal --cite-receipt <actions.get-receipt-id>
\end{lstlisting}

\noindent The \texttt{--cite-receipt} flag writes
\texttt{causal\_context.Scalar(\textit{prior\_invocation\_id})} into
the first step's invocation envelope, so the receipt log records
the chain
\texttt{actions.get $\to$ shop.cart\_contents $\to$ shop.checkout}.
A peer that fails to cite still gets a working run --- it just
loses the cross-state-object causal link, leaving each step's
receipt provable but the ``why this mission ran now'' answer
unrecorded.

Compare with the previous draft's static-directory design: the
same agent would fetch \texttt{checkout.eal},
\texttt{clear\_cart.eal}, \texttt{list\_items.eal}, an index
file, then filter --- four fetches plus a client-side parameter-
substitution pass for one intent, with no causal\_context tying
the run back to the page that prompted it.

\subsection{Hub mirror semantics}

\normsv{} The route action \texttt{actions:<verb>} translates an
HTTP request

\begin{lstlisting}
GET  /easynet-for-agents/<verb>?<query>
POST /easynet-for-agents/<verb>     {<json>}
\end{lstlisting}

into an EasyNet ability invocation
\texttt{<owner>.actions.<verb>} with the query string or JSON
body as args, identical to the \texttt{api:<verb>} dispatch.
The result is returned as JSON (or, when the result is a string
that the owner has flagged as \texttt{Content-Type:
text/eal}, served verbatim with that content type). Any verb
not in \texttt{actions\_surface\_hash} is rejected per WB-INV-2.

\begin{notebox}
The HTTP mirror's existence is convenience, not authority. A
peer EasyNet agent can equivalently call
\texttt{<owner>.actions.<verb>} directly through the bus.
Both paths produce the same kind of OperationalReceipt and the
same caller-context recording. The mirror exists primarily so
a browser-side script can build a UI on top of the actions
surface (e.g. a ``what can I do here'' dropdown) without
needing an EasyNet client.
\end{notebox}

\subsection{What the protocol does NOT enforce}

\derived{} Three deliberately-omitted enforcements:

\begin{itemize}[leftmargin=2em,itemsep=0.18em]
    \item \textbf{No cardinality between human\_view and
    actions.} An owner whose page shows ``12 items for sale''
    is not protocol-required to make
    \texttt{actions.list\_for\_page("/")} return 12 items.
    Honesty between views is an \emph{owner discipline}; the
    protocol provides the mechanism (both views are signed by
    the same owner; a peer auditing the chain can detect drift)
    but does not require equivalence as an invariant. v1's
    enforcement window is too narrow.
    \item \textbf{No standard EAL/mission text format version.}
    The body of \texttt{actions.get} is whatever the owner
    returns; the EAL grammar lives in its own RFC and evolves
    independently. Action descriptors carry \texttt{kind} so
    that a consumer can dispatch on body interpretation.
    \item \textbf{No authoritative ranking or relevance.}
    \texttt{actions.list\_for\_page} and \texttt{actions.search}
    return owner-defined rankings. A peer wanting independent
    ranking implements its own.
\end{itemize}

% ──────────────────────────────────────────────────────────────────────────
\section{Hub routing model}
\label{sec:hub-routing}

The Hub is bound to a single port (Phase 1: \texttt{127.0.0.1:8787}
by default). Every incoming HTTP request is processed in seven
steps:

\begin{enumerate}[leftmargin=2.2em,itemsep=0.16em]
    \item Mint envelope: \texttt{caller =
    principal/human-anon/<ip-hash>/<sid>};
    \texttt{caller\_context} from headers per main \S3.6.
    \item Parse path to \texttt{(realm, owner, slug, tail)}.
    \item Resolve state. If not Built, return per
    \S\ref{sec:hub-by-state}; if visibility forbids the
    caller, return 404.
    \item Match \texttt{tail} against the
    \texttt{routes\_manifest} rules in declared order.
    \item Dispatch on the matching action: \texttt{static} /
    \texttt{api} / \texttt{actions} / \texttt{ws} /
    \texttt{agent\_view} / \texttt{health} / \texttt{receipts} /
    \texttt{spa} fallback (see \S\ref{sec:hub-actions}).
    \item Build response (CSP, ETag, cache headers, content-type).
    \item Emit one OperationalReceipt with caller, caller\_context,
    matched rule, action, response status. Not canonical.
\end{enumerate}

\subsection{Hub behavior matrix per state}
\label{sec:hub-by-state}

\begin{center}
\renewcommand{\arraystretch}{1.25}
\begin{tabularx}{\linewidth}{l l X}
\toprule
\textbf{Canonical} & \textbf{Runtime} & \textbf{Hub response} \\
\midrule
\texttt{Created}  & *           & 503 ``not built''. \texttt{Retry-After: 60}. \\
\texttt{Building} & *           & 503 + small progress page (SSE of build.progress). \\
\texttt{Built}    & \texttt{Stopped}  & 503 ``not serving''. \\
\texttt{Built}    & \texttt{Starting} & 503 ``starting'', \texttt{Retry-After: 2}. \\
\texttt{Built}    & \texttt{Running}  & Full action dispatch (below). \\
\texttt{Built}    & \texttt{Crashed}  & 503 + reason hash; \texttt{Retry-After: 5}. \\
\texttt{Failed}   & *           & 503 ``last build failed'' + reason hash. \\
\texttt{Deleted}  & *           & 404. \\
\bottomrule
\end{tabularx}
\end{center}

\subsection{Action dispatch}
\label{sec:hub-actions}

\paragraph{static:<path>} TreeBlob path lookup, blob fetch,
extension-based MIME (allow-list), CSP from \texttt{csp\_hash},
immutable cache for \texttt{/assets/*}.

\paragraph{api:<verb>} Validate against
\texttt{api\_surface\_hash}, parse JSON body or query string
as args, invoke \texttt{<owner>.<verb>}, return result body
or structured error. Refusal: HTTP 404 (WB-INV-2).

\paragraph{actions:<verb>} Validate against
\texttt{actions\_surface\_hash}, parse JSON body or query
string as args, invoke \texttt{<owner>.actions.<verb>}, return
result body. Content-Type defaults to
\texttt{application/json}; if the result body declares
\texttt{Content-Type: text/eal}, the Hub forwards verbatim
under that type (so peers fetching an EAL file via the mirror
see proper text). Refusal: HTTP 404 (WB-INV-2).

\paragraph{ws:<path>} WebSocket upgrade; first text frame is a
\texttt{\{ ability, args \}} JSON; bind to ability stream;
forward subsequent frames.

\paragraph{agent\_view} Return JSON specified in
\S\ref{sec:agent-view}.

\paragraph{health} Return canonical and runtime states.

\paragraph{receipts} Return CanonicalReceipt history bounded by
\texttt{?from=...\&to=...}.

\paragraph{spa:<fallback-path>} Same as
\texttt{static:<fallback-path>} but \texttt{Cache-Control:
no-cache}.

% ──────────────────────────────────────────────────────────────────────────
\section{agent\_view}
\label{sec:agent-view}

\normsv{} \texttt{agent\_view} is structured \emph{metadata}
about the application: state, version, history, the surfaces
(api + actions) the application advertises. It is \emph{not}
the action library --- the actions live in the
\texttt{actions.*} namespace and are queried via that namespace,
not through agent\_view.

\begin{lstlisting}
GET /<realm>/<owner>/<slug>/__easynet/agent

200 OK
Content-Type: application/json

{
  "state_type":  "easynet.web_app",
  "state_key":   { "realm":"default","owner":"alice","slug":"shop" },
  "ura":         "easynet:///r/default/agent/alice/web_app/shop",

  "canonical": {
    "state":              "Built",
    "manifest":           { "name":"shop",
                            "description":"Tiny ceramic shop, agent-built.",
                            "visibility":"public" },
    "build_artifact_hash":"0x4d28ee17...",
    "version":            7,
    "etag":               "TSjuFw0pAvjpAQ=="
  },

  "runtime": {
    "state":     "Running",
    "started_at":"2026-04-28T09:11:02Z"
  },

  "ability_surface": [
    "webapp.update_config@v1",
    "webapp.build@v1",
    "webapp.stop@v1",
    "webapp.delete@v1"
  ],

  "api": {
    "endpoints": [
      { "verb":"shop.list_items",
        "http":"GET  /api/list_items",
        "schema_url":"/__easynet/schema/shop.list_items" },
      { "verb":"shop.checkout",
        "http":"POST /api/checkout",
        "schema_url":"/__easynet/schema/shop.checkout" }
    ]
  },

  "actions": {
    "namespace": "actions",
    "registered_verbs": [
      "actions.list_for_page",
      "actions.get",
      "actions.search"
    ],
    "http_mirror_prefix": "/easynet-for-agents/",
    "discovery_hint":
      "POST /easynet-for-agents/list_for_page {\"page_path\":\"<path>\"}"
  },

  "history": {
    "versions": 7,
    "latest_canonical_receipt_id": "r_c3d4...",
    "receipt_log_ura":
       "easynet:///r/default/agent/alice/web_app/shop/receipts"
  }
}
\end{lstlisting}

\textbf{Why agent\_view does not embed the action library.}
Embedding actions would force agent\_view to enumerate every
EAL fragment the owner ever wrote --- the same bulk-download
inefficiency that motivated the actions namespace in the first
place. agent\_view points at the namespace; the peer agent then
calls into the namespace for what it actually needs. The two
documents are complementary: agent\_view answers ``what is this
application,'' actions answers ``what can I do on it now.''

% ──────────────────────────────────────────────────────────────────────────
\section{Architecture}
\label{sec:architecture}

\subsection{Module map (EasyNet-Cli)}

\begin{lstlisting}
src/runtime/state/                    [shared with pages, P-A1]

src/runtime/blob_store/               [shared with pages, P-A1]
  tree.rs                             ← NEW: TreeBlob put/get,
                                       per-path lookup

src/runtime/agents/webapp_ability.rs  ← NEW (P-B2)
                                      register: create / build /
                                      update_config / serve / stop /
                                      delete / get / list

src/runtime/build_executor/           ← NEW (P-B1)
  mod.rs                              long-running canonical
                                      transition driver
  source_snapshot.rs                  Merkle hash of project dir
  spawn.rs                            subprocess + log streaming
  artifact_capture.rs                 dist tree -> TreeBlob

src/runtime/views/webapp/             ← NEW (P-B5)
  agent_view.rs                       JSON projection of §7
  health.rs
  receipts.rs

src/runtime/hub/                      [extends pages Hub from P-A4]
  router.rs                           ← extended for routes_manifest
                                       and the new actions: action
  static_handler.rs                   ← NEW (P-B3): TreeBlob fetch
  api_handler.rs                      ← NEW (P-B3): /api/<verb>
  actions_handler.rs                  ← NEW (P-B3): /easynet-for-agents/<verb>
  ws_handler.rs                       ← NEW (P-B4): /ws/* upgrade

src/runtime/skills/easynet-pages-author/  ← NEW (P-B5)
  SKILL.md                            agent-facing tutorial seeded
                                      into every workspace
\end{lstlisting}

\subsection{Daemon wiring}

The daemon at boot, on top of the pages wiring:

\begin{enumerate}[leftmargin=2em,itemsep=0.15em]
    \item Opens the build executor's working directory at
    \texttt{\textasciitilde/.easynet/build-cache/}.
    \item Registers \texttt{webapp\_ability} for each agent.
    \item For every web\_app whose canonical state is
    \texttt{Building} (i.e. a daemon-restart interrupted a
    build): marks Failed with reason ``daemon restart during
    build.''
    \item Seeds \texttt{easynet-pages-author} skill into each
    agent's workspace if not already present or if out of date.
\end{enumerate}

\subsection{Hub wiring}

The Hub binary (\texttt{easynet-hub}, from P-A4) gains in
P-B3/P-B4 the \texttt{static\_handler}, \texttt{api\_handler},
\texttt{actions\_handler}, and \texttt{ws\_handler} alongside the
page handler. Top-level dispatch:

\begin{lstlisting}
fn handle(req: HttpRequest) -> HttpResponse {
    let envelope = build_envelope(&req);
    let (realm, owner, slug, tail) = parse_path(&req.uri)?;
    let kind = state_store.classify(realm, owner, slug)?;

    match kind {
        Kind::Page    => pages_dispatch(envelope, owner, slug, tail),
        Kind::WebApp  => webapp_dispatch(envelope, owner, slug, tail),
        Kind::Unknown => http_404(),
    }
}
\end{lstlisting}

% ──────────────────────────────────────────────────────────────────────────
\section{Implementation plan}
\label{sec:impl-plan}

Five PR slices, sequenceable. P-B1 builds on P-A1 (blob store
and state store from pages); P-B3/B4 build on P-A4 (basic Hub).

\subsection{P-B1: build executor + TreeBlob}

\textbf{New files.} \texttt{src/runtime/build\_executor/} (4
modules), \texttt{src/runtime/blob\_store/tree.rs}.

\textbf{Public API.}
\begin{lstlisting}
pub trait BuildExecutor: Send + Sync {
    fn execute(
        &self,
        project_dir: &Path,
        config: &BuildConfig,
        progress: ProgressSink,
    ) -> Result<BuildOutcome>;
}
pub fn source_snapshot_hash(project_dir: &Path, ignore: &[Pattern])
    -> Result<[u8; 32]>;
pub fn capture_dist(dist_dir: &Path, blob: &dyn BlobStore)
    -> Result<(TreeBlob, [u8; 32])>;
\end{lstlisting}

\textbf{Tests.} Snapshot determinism, build smoke,
daemon-restart-during-build recovery.

\subsection{P-B2: webapp.* abilities}

\textbf{New files.} \texttt{src/runtime/agents/webapp\_ability.rs}.

\textbf{Behaviour.} Every transition in \S\ref{sec:state-machine}
is implemented, including the long-running build (per TR-INV-1).

\textbf{Tests.} Round-trip every transition; reject
\texttt{webapp.serve} from \texttt{Created}; optimistic
concurrency on \texttt{webapp.update\_config}; build progress
receipts loseable.

\subsection{P-B3: Hub static + API + actions routing}

\textbf{New files.} \texttt{src/runtime/hub/static\_handler.rs},
\texttt{api\_handler.rs}, \texttt{actions\_handler.rs}; extend
\texttt{router.rs} for \texttt{routes\_manifest}.

\textbf{Behaviour.}
\begin{enumerate}[leftmargin=2em,itemsep=0.1em]
    \item Static handler: TreeBlob path lookup, blob fetch,
    immutable cache for \texttt{/assets/*}.
    \item API handler: capture verb, validate against
    \texttt{api\_surface}, parse args, invoke
    \texttt{<owner>.<verb>}, map result.
    \item Actions handler: capture verb, validate against
    \texttt{actions\_surface}, parse args, invoke
    \texttt{<owner>.actions.<verb>}, map result. Defaults to
    \texttt{application/json}; respects
    \texttt{Content-Type: text/eal} returned by the ability.
\end{enumerate}

\textbf{Tests.} Integration: spin up a fixture web\_app whose
\texttt{api}.abilities contains \texttt{shop.echo} and whose
\texttt{actions}.abilities contains
\texttt{actions.list\_for\_page}; \texttt{POST /api/echo}
returns the echoed args; \texttt{POST /easynet-for-agents/list\_for\_page}
returns the action descriptors. Unauthorised verbs return 404.

\subsection{P-B4: WebSocket proxy + SPA fallback}

\textbf{New files.} \texttt{src/runtime/hub/ws\_handler.rs}.

\textbf{Behaviour.} WS upgrade on \texttt{/ws/*}; first text
message must be \texttt{\{ability,args\}} JSON; bind to
ability stream; forward frames. SPA fallback: serve
\texttt{dist/index.html} with \texttt{no-cache}.

\textbf{Tests.} WS streaming smoke, SPA fallback for unknown
paths.

\subsection{P-B5: agent\_view + meta endpoints + skill seed}

\textbf{New files.} \texttt{src/runtime/views/webapp/} (3
modules); minor router extensions for the
\texttt{/\_\_easynet/*} prefix; bundled skill at
\texttt{src/runtime/skills/easynet-pages-author/SKILL.md}.

\textbf{Behaviour.}
\begin{enumerate}[leftmargin=2em,itemsep=0.1em]
    \item \texttt{/\_\_easynet/agent} returns the JSON of \S7.
    \item \texttt{/\_\_easynet/health} returns canonical +
    runtime states.
    \item \texttt{/\_\_easynet/receipts} returns CanonicalReceipt
    history.
    \item \texttt{/\_\_easynet/schema/<verb>} returns the
    ability's input/output schemas.
    \item Skill seeder: at boot, daemon writes the bundled
    \texttt{easynet-pages-author} skill into each agent's
    workspace under \texttt{<workspace>/skills/} and
    \texttt{<workspace>/.agents/skills/}. Idempotent: skips
    when the skill is already present and up to date.
\end{enumerate}

\textbf{Tests.} agent\_view byte-determinism; skill seeder
idempotency; replay reconstructs canonical state from receipts.

% ──────────────────────────────────────────────────────────────────────────
\section{End-to-end worked example}
\label{sec:e2e}

\subsection{Setup}

\begin{lstlisting}
$ easynet agent add alice --type claude-code
$ easynet-daemon &        # seeds easynet-pages-author skill into
                          # alice's workspace at first contact
$ easynet-hub --bind 127.0.0.1:8787 &
\end{lstlisting}

\subsection{Author the project}

Claude Code, having read \texttt{easynet-pages-author}'s
SKILL.md, asks the daemon to mint state and create the project
directory:

\begin{lstlisting}
$ easynet ability invoke alice.webapp.create --args '{
    "slug":"shop",
    "manifest":{ "name":"shop",
                 "description":"Tiny ceramic shop, agent-built.",
                 "visibility":"public" }
}'
{ "state":"Created", "version":1,
  "project_dir":"~/.easynet/pages-project/shop" }
\end{lstlisting}

\noindent The agent then \texttt{cd}s into the project directory
and uses its toolchain:

\begin{lstlisting}
$ cd ~/.easynet/pages-project/shop
$ npm create vite@latest . -- --template react-ts
$ # ... write src/App.tsx, register API and actions abilities ...
$ # ... fill in easynet.app.toml as in §3.2 ...
$ easynet ability deploy abilities/shop.list_items     --node alice
$ easynet ability deploy abilities/shop.checkout       --node alice
$ easynet ability deploy abilities/actions.list_for_page --node alice
$ easynet ability deploy abilities/actions.get         --node alice
\end{lstlisting}

\subsection{Build and serve}

\begin{lstlisting}
$ easynet ability invoke alice.webapp.build --args '{ "slug":"shop" }' \
    --stream
[build.started   ] command="npm install && npm run build"
[build.progress  ] log_hash=0x12... 8% installed
[build.progress  ] log_hash=0x13... 47% bundled
[build.completed ] artifact_hash=0x4d28ee17... duration=38.2s
{ "state":"Built", "version":2,
  "build_artifact_hash":"0x4d28ee17..." }

$ easynet ability invoke alice.webapp.serve --args '{ "slug":"shop" }'
{ "runtime_state":"Running" }
\end{lstlisting}

\subsection{Browser fetch}

\begin{lstlisting}
$ curl -i http://127.0.0.1:8787/default/alice/shop/

HTTP/1.1 200 OK
Content-Type: text/html
ETag: "TSjuFw0pAvjpAQ=="
Cache-Control: max-age=60, must-revalidate
Content-Security-Policy: default-src 'self'; script-src 'self'; ...

<!doctype html>...
\end{lstlisting}

\subsection{API call from the browser}

\begin{lstlisting}
POST /default/alice/shop/api/checkout HTTP/1.1
Cookie: easynet_session=sess-abc
{ "item_id":"mug-12oz", "qty":1 }

→ HTTP/1.1 200 OK
{ "order_id":"ord_77a3f", "eta":"2026-04-29" }
\end{lstlisting}

\subsection{Peer agent discovers the actions}

A peer agent (bob) lands on the cart page and asks the owner
what it can do:

\begin{lstlisting}
$ curl -s -X POST http://127.0.0.1:8787/default/alice/shop/easynet-for-agents/list_for_page \
       -d '{"page_path":"/cart"}' | jq
[
  { "name":"checkout", "kind":"mission",
    "summary":"Buy the items currently in this user's cart.",
    "url":"/easynet-for-agents/get?name=checkout",
    "params_schema":{
      "address":{"type":"string","required":true}
    } },
  { "name":"clear_cart", "kind":"eal",
    "summary":"Empty the cart.",
    "url":"/easynet-for-agents/get?name=clear_cart" }
]

# bob asks the owner to specialise the mission for its concrete
# address input. Owner inlines the literal into a real EAL DAG
# (no $variable — EAL has none; see §6.4).
$ curl -s -X POST "http://127.0.0.1:8787/default/alice/shop/easynet-for-agents/get" \
       -d '{"name":"checkout",
            "params":{"address":"123 Main St"}}'

mission "checkout" {
  let cart = call "shop.cart_contents" on "alice" with {} timeout 5

  let order = call "shop.checkout" on "alice" with {
    items   = cart.output,
    address = "123 Main St"
  } timeout 30
}
\end{lstlisting}

\noindent Two RPCs, exact-match action, address inlined as a
literal by the owner. Bob now runs the returned text directly,
citing the \texttt{actions.get} receipt so the run's first step
carries the cross-state-object causal link:

\begin{lstlisting}
$ easynet mission run /tmp/checkout.eal --cite-receipt <actions.get-receipt-id>
\end{lstlisting}

Each \texttt{call} produces a CanonicalReceipt through the
existing dispatcher; the first step's invocation envelope
carries \texttt{causal\_context.Scalar} pointing at the
\texttt{actions.get} receipt, so an auditor walking the chain
sees \texttt{actions.list\_for\_page → actions.get →
shop.cart\_contents → shop.checkout} as one continuous causal
thread.

\subsection{Audit replay}

The web\_app's history shows three CanonicalReceipts (create,
build, build); each carries a \texttt{post\_state\_hash} and an
\texttt{owner\_signature}. The browser-driven checkout emitted a
CanonicalReceipt for an entirely different state\_type (an
\texttt{order} state object owned by alice or alice's order
service); that receipt also chains back to alice's signing key.
A peer reading both chains can reconstruct, byte-for-byte:

\begin{itemize}[leftmargin=2em,itemsep=0.12em]
    \item which build artifact the page was at when checkout
    happened;
    \item which API and actions endpoints were declared at
    that build;
    \item which mission body \texttt{actions.get} returned for
    the request that ultimately drove the checkout;
    \item that all of the above was signed by alice.
\end{itemize}

\noindent This is the audit story decisively unavailable on the
traditional web.

% ──────────────────────────────────────────────────────────────────────────
\section{Open questions deferred to RFC-006-B v2}
\label{sec:open}

\begin{enumerate}[leftmargin=2.2em,itemsep=0.15em]
    \item SSR with persistent process (\texttt{serve.mode = process}).
    \item Cross-node hubs, federated reads.
    \item Signed asset bundles for SRI.
    \item Build sandboxing.
    \item Multi-version routing
    (\texttt{webapp.tag} transitions).
    \item dev-mode reverse-proxy
    (\texttt{serve.mode = dev-proxy}) for HMR / Vite dev server.
\end{enumerate}

\subsection{Coverage check}

\begin{center}
\renewcommand{\arraystretch}{1.2}
\begin{tabularx}{\linewidth}{l X l}
\toprule
\textbf{Main rule} & \textbf{Discharged by} & \textbf{Where} \\
\midrule
SO-INV-1 (key immutability)              & state\_key parser refuses key updates                            & \S\ref{sec:state-machine} \\
CSH-INV-1 (hash over projection)         & schema-driven $\pi$, JCS, SHA-256                                 & \S\ref{sec:state-machine} \\
CSH-INV-2 (content as blob)              & TreeBlob + per-file blobs; canonical state holds only hashes      & \S\ref{sec:build-pipeline} \\
RC-INV-1 (receipt binding)               & every webapp receipt carries (state\_key, transition\_id, attempt\_id)  & \S\ref{sec:state-machine} \\
RC-INV-2 (optimistic concurrency)        & ETag check on update\_config                                       & P-B2 \\
TA-INV-1 (transition\_id required)       & every \texttt{webapp.<verb>@v1} validated                          & \S\ref{sec:state-machine} \\
VP-INV-1 (ability\_surface from state)   & \texttt{ability\_surface} derived from state\_value                & \S\ref{sec:agent-view} \\
TR-INV-1 (long-running terminal)         & build emits exactly one terminal CanonicalReceipt                  & \S\ref{sec:build-pipeline} \\
TR-INV-2 (progress not canonical)        & \texttt{build.progress} are OperationalReceipts                    & \S\ref{sec:build-pipeline} \\
TR-INV-5 (runtime overlay independent)   & runtime crash leaves canonical untouched                           & \S\ref{sec:state-machine} \\
TR-INV-7 (multi-view first-class)        & human (HTML+API+WS) and agent\_view both implemented               & \S\ref{sec:hub-routing}, \S\ref{sec:agent-view} \\
TR-INV-8 (cross-state-space matrix)      & canonical $\times$ runtime admissibility enforced                  & \S\ref{sec:state-machine} \\
TR-INV-11 (URA-receipt traceability)     & \texttt{/\_\_easynet/receipts} + state\_key index                  & \S\ref{sec:e2e} \\
TR-INV-12 (principal/human-anon)         & Hub mints the URI on every request                                 & \S\ref{sec:hub-routing} \\
TR-INV-13 (caller\_context)              & populated from headers; recorded in operational receipts          & \S\ref{sec:hub-routing}, \S\ref{sec:e2e} \\
WB-INV-1 (artifact-from-blob-only)       & Hub never reads project's dist/                                    & \S\ref{sec:project} \\
WB-INV-2 (api/actions surface enforcement)& Hub refuses unlisted verbs with 404                                & \S\ref{sec:project} \\
\bottomrule
\end{tabularx}
\end{center}

\paragraph{Closing.} The agent-native web's discoverability
problem is now solved at the right layer. A peer agent landing
on a page asks \texttt{actions.list\_for\_page} and gets exactly
the EAL/missions relevant to that page, parameterised by its own
intent --- not a sitemap to filter, not a manifest to download.
Authoring is a learnable skill, not folklore: every agent
acquires \texttt{easynet-pages-author} on first contact. The
project lives at a stable, agent-independent location
(\texttt{\textasciitilde/.easynet/pages-project/}) so the same
codebase reads the same way regardless of which agent operates
it. Implementing P-B1 through P-B5 closes Phase 1.

\vfill
\hrule
\vspace{0.4em}
\noindent\small
\textit{The web evolved from documents to applications;
agent-native web evolves from passive views to queryable
actions. Pages first, web\_apps next, actions are how peer
agents make use of either without bulk-downloading anything.}

\end{document}
